\documentclass[a4paper, 14pt]{extarticle}

% Поля
%--------------------------------------
\usepackage{geometry}
\geometry{a4paper,tmargin=2cm,bmargin=2cm,lmargin=3cm,rmargin=1cm}
%--------------------------------------


%Russian-specific packages
%--------------------------------------
\usepackage[T2A]{fontenc}
\usepackage[utf8]{inputenc} 
\usepackage[english, main=russian]{babel}
%--------------------------------------

\usepackage{textcomp}

% Красная строка
%--------------------------------------
\usepackage{indentfirst}               
%--------------------------------------             


%Graphics
%--------------------------------------
\usepackage{graphicx}
\graphicspath{ {./images/} }
\usepackage{wrapfig}
%--------------------------------------

% Полуторный интервал
%--------------------------------------
\linespread{1.3}                    
%--------------------------------------

%Выравнивание и переносы
%--------------------------------------
% Избавляемся от переполнений
\sloppy
% Запрещаем разрыв страницы после первой строки абзаца
\clubpenalty=10000
% Запрещаем разрыв страницы после последней строки абзаца
\widowpenalty=10000
%--------------------------------------

%Списки
\usepackage{enumitem}

%Подписи
\usepackage{caption} 

%Гиперссылки
\usepackage{hyperref}

\hypersetup {
	unicode=true
}

%Рисунки
%--------------------------------------
\DeclareCaptionLabelSeparator*{emdash}{~--- }
\captionsetup[figure]{labelsep=emdash,font=onehalfspacing,position=bottom}
%--------------------------------------

\usepackage{tempora}
\usepackage{amsmath}
\usepackage{color}
\usepackage{listings}
\lstset{
  belowcaptionskip=1\baselineskip,
  breaklines=true,
  frame=L,
  xleftmargin=\parindent,
  language=Python,
  showstringspaces=false,
  basicstyle=\footnotesize\ttfamily,
  keywordstyle=\bfseries\color{blue},
  commentstyle=\itshape\color{purple},
  identifierstyle=\color{black},
  stringstyle=\color{red},
}

%--------------------------------------
%			НАЧАЛО ДОКУМЕНТА
%--------------------------------------

\begin{document}

%--------------------------------------
%			ТИТУЛЬНЫЙ ЛИСТ
%--------------------------------------
\begin{titlepage}
\thispagestyle{empty}
\newpage


%Шапка титульного листа
%--------------------------------------
\vspace*{-60pt}
\hspace{-65pt}
\begin{minipage}{0.3\textwidth}
\hspace*{-20pt}\centering

\includegraphics[width=\textwidth]{emblem}
\end{minipage}
\begin{minipage}{0.67\textwidth}\small \textbf{
\vspace*{-0.7ex}
\hspace*{-6pt}\centerline{Министерство науки и высшего образования Российской Федерации}
\vspace*{-0.7ex}
\centerline{Федеральное государственное бюджетное образовательное учреждение }
\vspace*{-0.7ex}
\centerline{высшего образования}
\vspace*{-0.7ex}
\centerline{<<Московский государственный технический университет}
\vspace*{-0.7ex}
\centerline{имени Н.Э. Баумана}
\vspace*{-0.7ex}
\centerline{(национальный исследовательский университет)>>}
\vspace*{-0.7ex}
\centerline{(МГТУ им. Н.Э. Баумана)}}
\end{minipage}
%--------------------------------------

%Полосы
%--------------------------------------
\vspace{-25pt}
\hspace{-35pt}\rule{\textwidth}{2.3pt}

\vspace*{-20.3pt}
\hspace{-35pt}\rule{\textwidth}{0.4pt}
%--------------------------------------

\vspace{1.5ex}
\hspace{-35pt} \noindent \small ФАКУЛЬТЕТ\hspace{80pt} <<Информатика и системы управления>>

\vspace*{-16pt}
\hspace{47pt}\rule{0.83\textwidth}{0.4pt}

\vspace{0.5ex}
\hspace{-35pt} \noindent \small КАФЕДРА\hspace{50pt} <<Теоретическая информатика и компьютерные технологии>>

\vspace*{-16pt}
\hspace{30pt}\rule{0.866\textwidth}{0.4pt}
  
\vspace{11em}

\begin{center}
\Large {\bf Лабораторная работа № 3} \\ 
\large {\bf<<Оценка ошибки округления и ошибки в передаваемых данных>>} \\ 
\large {\bf по курсу <<Численные методы линейной алгебры>>} \\ 
\end{center}\normalsize

\vspace{8em}


\begin{flushright}
  {Студент группы ИУ9-71Б Лысенко В. В.\hspace*{15pt} \\
  \vspace{2ex}
  Преподаватель Посевин Д. П.\hspace*{15pt}}
\end{flushright}

\bigskip

\vfill
 

\begin{center}
\textsl{Москва 2024}
\end{center}
\end{titlepage}
%--------------------------------------
%		КОНЕЦ ТИТУЛЬНОГО ЛИСТА
%--------------------------------------

\renewcommand{\ttdefault}{pcr}

\setlength{\tabcolsep}{3pt}
\newpage
\setcounter{page}{2}

\section{Цель работы}\label{Sect::task}
\par
\renewcommand{\labelenumii}{\arabic{enumi}.\arabic{enumii}.}
Реализовать автоматическое вычисление оценки относительной ошибки округления и ошибки в передаваемых данных.

\section{Задачи}\label{Sect::task}
\par
\renewcommand{\labelenumii}{\arabic{enumi}.\arabic{enumii}.}
\begin{enumerate} 
    \item Реализовать автоматическое вычисление оценки относительной ошибки округления и ошибки в передаваемых данных для больших объёмов данных.
    \item Реализовать автоматическое вычисление оценки относительной ошибки округления и ошибки в передаваемых данных для значений из индивидуального варианта. 
    \item Сравнение вычисленных данных с полученными на бумаге.
\end{enumerate}

\pagebreak
\section{Практическая реализация}
\par
Для получения ошибки округления необходима функция, с погрешностью вычисляющая вектор х в выражении Ax = f. В качестве таковой была взята функция из предыдущей лабораторной работы, реализующая классический метод Гаусса. Данный метод был переопределён таким образом, чтобы возвращать помимо вектора х ещё и g(A) (коэффициент роста).
\par
При вычислении ошибки передаваемых данных для получения вектора х использовались уже имеющиеся возможности языка. 
\par
При выдаче задания было сказано находить норму по формуле Евклида, однако в письменной работе использовалась формула равномерной нормы, поэтому в данной лабораторной работе были реализованы методы для нахождения обеих норм ( norm() и matrix-norm() - для поиска норм вектора и матрицы по Евклиду, inf-norm() и inf-matrix-norm() - для поиска равномерных норм вектора и матрицы).
\par

Ниже представлен код программы:
\begin{lstlisting}
n = 2
p = 2.0
t = 26
cup = 100

function gauss_classic(A, F)
    max = maximum(abs.(A))
    max_elems = zeros(Float64, n-1)

    for k in 1:n-1
        for i in k+1:n
            factor = A[i, k] / A[k, k]
            A[i, k:end] .-= factor * A[k, k:end]
            F[i] -= factor * F[k]
        end

        max_elems[k] = maximum(abs.(A)) / max     
    end

    g_a = maximum(max_elems)

    X = zeros(Float64, n)
    for i in n:-1:1
        X[i] = F[i] / A[i, i]
        for j in i-1:-1:1
            F[j] -= A[j, i] * X[i]
        end
    end

    return g_a, X
end


function norm(x)
    sum = 0.0
    for element in x
        sum += element^2
    end
    return sqrt(sum)
end

function matrix_norm(A)
    sum = 0.0
    for i in 1:n
        for j in 1:n
            #print(A[i, j], " ")
            sum += A[i, j] ^ 2
        end
        #println()
    end
    return sqrt(sum)
end

function inf_norm(x)
    return maximum(abs.(x))
end

function inf_matrix_norm(A)
    return maximum([sum(abs.(A[i, :])) for i in 1:size(A, 1)])
end



#X = ones(Float64, n)
#A = rand(-cup:cup, n, n)
#A = Float64.(A)
#F = [Float64(x) for x in (A * X)]

X = [2.625 1.5]
A = [[48.0, 340.0] [340.0, 23.0]]
F = [636.0 927.0]


println("---------------------------------------------------")

g_A, X0 = gauss_classic(copy(A), copy(F))


delta_X = vec(X) - X0

norm_X = inf_norm(X)
norm_delta_X = inf_norm(delta_X)

error = norm_delta_X / norm_X

nu_A = inf_matrix_norm(A) * inf_matrix_norm(inv(A))

#println("X = ", X)
#println("X0 = ", X0)
#println("delta_X = ", delta_X)
println("norm(delta_X) = ", norm_delta_X)

println()

println("nu(A) = ", nu_A)
println("n = ", n)
println("g(A) = ", g_A)
println("p^(-t) = ", p^(-t))

rounding_error = nu_A * n * g_A * (p^(Float64(-t)))

println("\n", error,  " <= ", rounding_error, "\t<--- rounding_error")
println("---------------------------------------------------")

A1 = copy(A)
F1 = copy(F)
X1 = copy(X)

delta_A1 = zeros(Float64, n, n)
delta_F1 = zeros(Float64, n)
#delta_X1 = zeros(Float64, n)

delta_A1 .+= 25
delta_F1 .+= 2.6

wrong_A1 = A1 + delta_A1
wrong_F1 = vec(F1) + vec(delta_F1)


X01 = wrong_A1 \ wrong_F1
#println(X01)

delta_X1 = vec(X1) - vec(X01)
#println(delta_X1)
nu_A1 = inf_matrix_norm(A1) * inf_matrix_norm(inv(A1))

error1 = inf_norm(delta_X1) / inf_norm(X1)

data_error = nu_A1 * ( inf_norm(delta_F1) / inf_norm(F1) + inf_matrix_norm(delta_A1) / inf_matrix_norm(A1))
println("\n", error1,  " <= ", data_error, "\t<--- data_error")
\end{lstlisting}

\section{Результаты}
\par
Результат с данными из индивидуального варианта:
\begin{center}
\includegraphics[scale=0.75]{g1}
\par
Как видно, результаты совпали с вычисленными на бумаге.
\end{center}
\par

Результат с большими данными (размерность матрицы и векторов 100):
\begin{center}
\includegraphics[scale=0.75]{g2}
\par
\end{center}

\section{Выводы}
В ходе выполнения лаборатороной работы было реализовано автоматическое вычисление оценки относительной ошибки округления и ошибки в передаваемых данных для больших объёмов данных. Корректность работы программы была так же протестирована на данных из индивидуального варианта, в результате чего полученный результат сошёлся с ожидаемым.
\end{document}
